%! Author = alan
%! Date = 15/03/2025

% Preamble
\documentclass[a4paper, 11pt]{article}
\usepackage[left=3cm, top=3cm, right=2cm, bottom=2cm]{geometry} %pacote que define as margens
\usepackage[scaled]{helvet}
\renewcommand\familydefault{\sfdefault} % define a fonte ARIAL como padrão do texto
\usepackage{graphicx} %pacote que permite a inserção de imagens
\usepackage{setspace} %pacote que permite especificar o espaçamento entre linhas no documento
\usepackage{indentfirst}
\setlength{\parindent}{1cm} %pacote que faz indentação na primeira linha do parágrafo
\usepackage{float} %Pacote que permite tabelas flutuarem em qualquer posição
\usepackage[alf,abnt-repeated-title-omit=yes,abnt-emphasize=bf,abnt-etal-list=0]{abntex2cite} %define os parâmetros necessários para as referências e citações de acordo com a ABNT
\usepackage{amsthm} % Pacote que permite a utilização de teoremas
\usepackage{polyglossia}
\usepackage{amsfonts} % Pacote que permite a utilização de símbolos matemáticos
\usepackage{caption} % Pacote que permite a inserção de legendas em figuras
\usepackage{subcaption} % Pacote que permite que imagens possam ser inseridas em composições de imagens lado a lado com o \subfigure
\usepackage{amssymb} %4 Pacotes responsáveis pela formatação matemática
\usepackage{rotating} %permite que textos, figuras e tabelas possam ser rotacionados
\usepackage{array}    % permite e vetorização de tabelas
\usepackage{longtable} % permite a elaboração de tabelas complexas
\usepackage{colortbl} % permite a utilização de paletas de cores
\usepackage{tikz} % criação de objetos vetoriais
\usetikzlibrary{shapes, arrows.meta, positioning} % permite a utilização de formas e setas
\usepackage{mathtools} % permite a utilização de equações matemáticas
\usepackage{pdflscape} % permite alterar o estilo da pagina entre portrato e paisagem
\usepackage{sectsty} % permite a alteração do estilo das seções
\usepackage{url} % permite a utilização de links
\usepackage{pgfplots} % permite a criação de gráficos
\usepackage[table,xcdraw]{xcolor} % permite a utilização de cores
\usepackage[brazilian,hyperpageref]{backref}    % Paginas com as citações na bibl
\usepackage{hyperref}
\usepackage{enumitem}
\usepackage{microtype} % para melhorias de justificação
\hypersetup{
    colorlinks=true,
    linkcolor=black, % alterado para black, antes darkblue por bug
    filecolor=magenta,
    urlcolor=blue,
    citecolor=black,
    pdftitle={mybib},
    pdfpagemode=FullScreen
}
\usepackage{pgfgantt}  % gera gantt em latex
%\usepackage[acronym]{glossaries} % cria glossário de siglas

\usepackage[most]{tcolorbox}
\usepackage{listings}
\usepackage{cite}
\usepackage{mdwtab}
\usepackage{natbib}
\pagestyle{myheadings} %pacote que define o cabeçalho com o número de página no canto superior direito
\usepackage[utf8]{inputenc} %pacote que permite a utilização de acentos
%\usepackage[brazil]{babel} %pacote que define o idioma
\setmainlanguage{brazil}

\sectionfont{\fontsize{12}{12}\selectfont}
\subsectionfont{\fontsize{12}{12}\selectfont}
\subsubsectionfont{\fontsize{12}{12}\selectfont}
\lstset{
    inputencoding=utf8,
    extendedchars=true,
    literate={á}{{\'a}}1 {é}{{\'e}}1 {í}{{\'i}}1 {ó}{{\'o}}1 {ú}{{\'u}}1
        {Á}{{\'A}}1 {É}{{\'E}}1 {Í}{{\'I}}1 {Ó}{{\'O}}1 {Ú}{{\'U}}1
        {à}{{\`a}}1 {è}{{\`e}}1 {ì}{{\`i}}1 {ò}{{\`o}}1 {ù}{{\`u}}1
        {À}{{\`A}}1 {È}{{\`E}}1 {Ì}{{\`I}}1 {Ò}{{\`O}}1 {Ù}{{\`U}}1
        {ã}{{\~a}}1 {õ}{{\~o}}1 {Ã}{{\~A}}1 {Õ}{{\~O}}1
        {â}{{\^a}}1 {ê}{{\^e}}1 {î}{{\^i}}1 {ô}{{\^o}}1 {û}{{\^u}}1
        {Â}{{\^A}}1 {Ê}{{\^E}}1 {Î}{{\^I}}1 {Ô}{{\^O}}1 {Û}{{\^U}}1
        {ç}{{\c{c}}}1 {Ç}{{\c{C}}}1
}

\definecolor{codegreen}{rgb}{0,0.6,0}
\definecolor{codegray}{rgb}{0.5,0.5,0.5}
\definecolor{codepurple}{rgb}{0.58,0,0.82}
\definecolor{backcolour}{rgb}{0.95,0.95,0.92}

\lstdefinestyle{mystyle}{
    backgroundcolor=\color{backcolour},
    commentstyle=\color{codegreen},
    keywordstyle=\color{magenta},
    numberstyle=\tiny\color{codegray},
    stringstyle=\color{codepurple},
    basicstyle=\ttfamily\footnotesize,
    breakatwhitespace=false,
    breaklines=true,
    captionpos=b,
    keepspaces=true,
    numbers=left,
    numbersep=5pt,
    showspaces=false,
    showstringspaces=false,
    showtabs=false,
    tabsize=2
}

\lstset{style=mystyle}
\citeoption{abnt-etal-cite=0, abnt-and-type=&}


\begin{document}
    \thispagestyle{empty}

\begin{center}
	\begin{figure}[h]
  \centering
		\includegraphics[width=0.21\linewidth]{images/ufpa}
		\label{fig:ufpa}
	\end{figure}


	\vspace{1cm}
	\large \uppercase{UNIVERSIDADE FEDERAL DO PARÁ}\\
	\large \uppercase{INSTITUTO DE GEOCIÊNCIAS}\\
	\vspace{7cm}
	\large \uppercase{DESENVOLVIMENTO DE SENSOR DE SALINIDADE}\\
	\vspace{1cm}
	\large \uppercase {LABORATÓRIO DE OCEANOGRAFIA} \\
	\vspace{7cm}
	\large {BELÉM/PA \\ 2025}

 \newpage
 \thispagestyle{empty}
 \large \uppercase{alan henrique pereira miranda - 202102140072}\\
 \vspace{7cm}

 \singlespacing
 \hspace{8cm} % posicionando a caixa de texto
 \begin{minipage}{7cm}
	Atividade referênte ao desenvolvimento de projeto de pesquisa com o laboratório de oceanografia. \\

	Profa. Dr.: Renan Peixoto\\
	\vspace{1cm}

	Belém-PA 3 de fevereiro de 2025
	\vspace{4cm}
\end{minipage}

\onehalfspacing
\begin{center}

	EXAMINADOR\\
	\vspace{3cm}
	\rule{10cm}{0.15mm} \\
	Prof. Dr: Renan Peixoto\\
	Universidade Federal do Pará - UFPA
\end{center}
\newpage

\begin{center}
    % BLOCO DE FIGURAS
    \thispagestyle{empty}
	\listoffigures
	\newpage
    % SUMARIO
    \thispagestyle{empty}
    \tableofcontents

\end{center}

\newpage
\thispagestyle{empty}

\end{center}
    \section{Introdução}
\label{sec:introducao}

\subsection*{A Evolução dos Sensores: Uma Perspectiva Histórica}
Nos tempos atuais, os sensores são algo comum no cotidiano cujo o comportamento passa desapercebido, \cite{ferreira_2016_biblioteca} porém permeia diversos recursos e componentes que utilizamos no dia a dia. \citeonline{stenstrasser_2017_projeto, cengel_2012_transferencia, ferreira_2016_biblioteca, orlando_2014_aula}


    \bibliographystyle{abntex2-alf}
    \newpage
    \bibliography{main}
\end{document}